\chapter{Záver}
\label{sec:Zaver}

Táto práca dokumentuje priebeh mojej odbornej praxe vo firme Railsformers s.r.o. V~súlade so zadaním som postupne popísal odborné zameranie firmy, zoznam zadaných úloh, zvolený postup riešenia, uplatnené aj chýbajúce znalosti a~dosiahnuté výsledky s~celkovým zhodnotením.

Pracovné prostredie vo firme Railsformers bolo mimoriadne podporujúce. Ako malá firma s~priateľskou atmosférou mi poskytla priestor na samostatnú prácu, no zároveň som mal vždy k~dispozícii pomoc skúsenejších kolegov. Flexibilný prístup k~organizácii práce, možnosť home-office a~pravidelné konzultácie s~mentorom vytvorili ideálne podmienky pre učenie sa nových technológií. Ocenil som predovšetkým ochotu tímu vysvetľovať nielen technické aspekty, ale aj širší kontext rozhodnutí a~osvedčené postupy z~praxe.

Absolvovanie odbornej praxe mi prinieslo realistický pohľad na prácu v~softvérovej firme. Pred praxou som mal predstavu o~vývoji softvéru založenú prevažne na školských projektoch. Reálna práca však ukázala, že rovnako dôležité ako samotné programovanie sú aj komunikácia, dokumentácia, riešenie neočakávaných problémov a~schopnosť prispôsobiť sa meniacim požiadavkám. Táto skúsenosť mi dala sebavedomie vstúpiť do profesionálneho IT prostredia.

Budúcim študentom jednoznačne odporúčam absolvovať odbornú prax. Je to príležitosť overiť si teoretické znalosti v~praxi, získať cenné kontakty a~zistiť, ktorým oblastiam IT sa chcete venovať. Práca na reálnom projekte s~konkrétnymi deadlinmi a~požiadavkami je skúsenosť, ktorú žiadny školský predmet nedokáže plnohodnotne nahradiť.

\endinput
